\begin{lemma}[The \(C=39\) arithmetic front]\label{lem:n11-C39-front-v2}
Assume that eleven admissible points have \(39\) proper circles, contain a
proper five-circle, and have no proper circle larger than five.  Every
generalized block has size at most five, and the number \(L\) of maximal
lines belongs to \(\{12,15\}\).

Fix a proper five-circle \(\Gamma\), with outsider set \(X\).  Let \(A_{ij}\)
count nonselected blocks containing \(i\) points of \(\Gamma\) and \(j\)
outsiders, and put
\[
 H=A_{22}+3A_{23},\quad x=A_{23},\quad
 r=A_{13},\quad s=A_{14},\quad
 q=A_{04}+3A_{05},\quad y=A_{05}.
\]
If \(L=3k\), and if
\[
 B_0=A_{03}+A_{04}+A_{05},\qquad
 T_0=A_{03}+4A_{04}+10A_{05},
\]
then
\begin{align}
A_{22}&=H-3x,&A_{21}&=60-2H+3x,\notag\\
A_{12}&=75-2H-3r-6s,&
B_0&=3H+L-97-x+2r+5s,\notag\\
T_0&=20-x-r-4s,&
q&=39-H-k-r-3s. \label{eq:C39-full-front-v2}
\end{align}
Moreover
\begin{equation}\label{eq:C39-rich-cap-v2}
 0\le q\le3,\qquad y=1\Longrightarrow q=3,
\end{equation}
and, if \(e_\Gamma,e_X\) count points with \(d_3=9\) on the two sides,
\begin{equation}\label{eq:C39-high-split-v2}
 e_\Gamma=55-2H+2x-r-2s,\qquad
 e_X=2k-r-2s-2q+2y.
\end{equation}
Every point has \(d_3\in\{6,9\}\) and \(d_5\le4\).
\end{lemma}

\begin{proof}
Let \(A\) be a maximal \(a\)-line and put \(b=11-a\).  For each pair of
outsiders, its pair-hosts on \(A\) form a matching, so the line pencil gives
the uniform lower bound
\[
 C\ge b\binom a2-\binom b2\floor{a/2}.
\]
For \(a=6,7,8,9,10\) its values are respectively
\(45,66,72,68,45\), so a line of size at least six is impossible.  Let
\(N_j\) be the number of generalized \(j\)-blocks.  Counting blocks and
triples gives
\[
 N_3+N_4+N_5=39+L,\qquad N_3+4N_4+10N_5=165.
\]
Consequently
\begin{equation}\label{eq:C39-global-blocks-v2}
 L=3k,\qquad N_4+3N_5=42-k,\qquad
 N_3=-3+4k+2N_5.
\end{equation}

At a pivot, inversion gives
\[
 d_3+3d_4+6d_5=45.
\]
Kelly--Moser gives \(d_3\ge5\).  Deleting the pivot destroys at most
\(39-f(10)=6\) three-circles, and its three-lines use disjoint pairs of the
other ten points, so \(d_3\le11\).  The displayed congruence now yields
\(d_3\in\{6,9\}\).

If five four-point image lines existed at a pivot, their twenty incidences
on ten points and
\(\sum_u\binom{r_u}{2}\le\binom52\) would force every image point to be
the intersection of exactly two of the lines.  A further image line distinct
from those five contains at most two such pair-labelled points: three labels
would either share an index, giving one of the five lines, or be three
disjoint edges of \(K_5\), which is impossible.  The link equation nevertheless
requires a further three-point line.  Six four-point image lines are excluded
already by \(\sum_u\binom{r_u}{2}\ge18>\binom62\).  Hence \(d_5\le4\).

Melchior on the original lines gives
\[
 3\ell_3+7\ell_4+12\ell_5\le52,
\]
so \(L\le15\).  Summing the pivot slack and the deletion bound gives
\[
 \sum_p(d_3(p)-3-d_5(p))=4L+N_5-42\ge0,
 \qquad 2N_5\le25-k.
\]
These inequalities exclude \(L=0,3,6\).  For \(L=9\), put
\(E=\#\{p:d_3(p)=9\}=2N_5-13\).  Then \(7\le N_5\le11\).
Splitting the pivot slack between \(\Gamma\) and \(X\) gives
\[
 45+4y-4N_5\le3r+7s+6q\le39+4y-3N_5,
 \qquad r+3s+q\ge6.
\]
For \(q=0,1,2,3\), these intervals leave only
\[
 (N_5,r,s)=(7,6,0),(7,1,2),(8,0,2).
\]
The first and third rows have \(H=30\), hence no deficient outsider pair,
although respectively an \(A_{13}\) or an \(A_{14}\) block would require
one.  The middle row has \(H=29\), whose sole deficient edge cannot meet all
four outsider triples of an \(A_{14}\) block.  Thus \(L\ne9\), proving
\(L\in\{12,15\}\).

The three triple partitions relative to \(\Gamma\) give
\eqref{eq:C39-full-front-v2}.  Distinct outsider \(K_4\)'s have disjoint
complement pairs, while an outsider \(K_5\) excludes every \(K_4\).
Therefore either \(A_{05}=1,A_{04}=0\), or \(A_{05}=0,A_{04}\le3\), which
is \eqref{eq:C39-rich-cap-v2}.  Finally, summing three-block incidences on
\(\Gamma\) and on \(X\) gives \eqref{eq:C39-high-split-v2}.
\end{proof}

\begin{lemma}[Added-centre and singleton interfaces]
\label{lem:n11-C39-interfaces-v2}
Under the hypotheses of \cref{lem:n11-C39-front-v2}, put \(m=N_5\).  Then
\begin{equation}\label{eq:C39-line-cut-v2}
 L=12:\quad7\ell_4+16\ell_5\le m+8,
 \qquad
 L=15:\quad7\ell_4+16\ell_5\le m-7.
\end{equation}
If two generalized five-blocks meet in the single point \(p\), then
\begin{equation}\label{eq:C39-singleton-v2}
 (d_3(p),d_5(p))=(9,2).
\end{equation}
Moreover, every generalized five-block contains at most two points of type
\((d_3,d_4,d_5)=(9,4,4)\).
\end{lemma}

\begin{proof}
Write \(\sigma_p=d_3(p)-3-d_5(p)\).  The inversion images are noncollinear:
otherwise inversion back would put all eleven original points on one
generalized block.  Thus \cref{prop:universal-melchior-rows} applies,
gives \(\sigma_p\ge0\), and its restored-centre row is
\[
 3\ell_3(p)+\sum_{j\ge4}j\ell_j(p)\le10+\sigma_p.
\]
After summing over the eleven pivots and using
\eqref{eq:C39-global-blocks-v2},
\[
 9\ell_3+16\ell_4+25\ell_5
 \le110+\sum_p\sigma_p,
 \qquad \sum_p\sigma_p=4L+m-42.
\]
Since \(L=\ell_3+\ell_4+\ell_5\), this is exactly
\eqref{eq:C39-line-cut-v2}.

For the singleton assertion, invert at its carrier \(p\).  The two
five-blocks become disjoint four-point link lines, so \(z=d_5(p)\ge2\).
The cap \(z\le4\), pointwise \(\sigma_p\ge0\), and
\(d_3\in\{6,9\}\) leave only
\[
 (d_3,z)=(6,2),(6,3),(9,2),(9,3),(9,4).
\]
The two cases \(z=3\) are excluded by \cref{lem:n11-K3}(K3.1).  In the
case \((d_3,z)=(9,4)\), the link equation gives \(d_4=4\), so the
complete-quadrangle conclusion (K3.2) says that every pair of base lines
meets at a selected link point.

It remains to exclude \((d_3,z)=(6,2)\).  Dualizing the ten-point link turns
the two disjoint four-point lines into two fourfold points.  The ten dual
lines split into families of sizes \(4,4,2\); after removing pairs already
meeting at the two fourfold points, the four pair classes have sizes
\[
 16,\qquad8,\qquad8,\qquad1.
\]
The link equation requires nine triple vertices.  If the sole pair in the
last class is unused, all nine triples require one pair from each eight-edge
side class.  If it is used, the other eight triples exhaust both side
classes before the exceptional triple uses two further side edges.  Both
alternatives are impossible.  This proves \eqref{eq:C39-singleton-v2} for
circles and for the possible five-line alike.

Finally, a point of type \((9,4,4)\) supplies a harmonic complement on every
incident five-block by (K3.2).  The real cross-ratio argument (K3.3) allows
at most two such complements on five points, proving the last assertion.
\end{proof}

\begin{lemma}[The \(L=15\) harmonic collapse]
\label{lem:n11-C39-L15-v2}
The \(C=39,L=15\) branch is impossible.
\end{lemma}

\begin{proof}
Let \(z_p=d_5(p)\), and call \(p\) high when \(d_3(p)=9\).  From
\eqref{eq:C39-global-blocks-v2},
\[
 E:=\#\{p:p\text{ is high}\}=2m-5,
 \qquad \sum_pz_p=5m.
\]
Thus \(3\le m\le8\).  The second cut in
\eqref{eq:C39-line-cut-v2} forces \(m\in\{7,8\}\) and
\(\ell_4=\ell_5=0\); in particular all fifteen maximal lines are triples
and all generalized five-blocks are proper circles.

Define the nonnegative integral restored-centre slack
\[
 \kappa_p=10+\sigma_p-3\ell_3(p).
\]
Summing gives \(\sum_p\kappa_p=m-7\), while pointwise
\begin{equation}\label{eq:C39-kappa-point-v2}
 3\ell_3(p)=
 \begin{cases}
 13-z_p-\kappa_p,&p\text{ low},\\
 16-z_p-\kappa_p,&p\text{ high}.
 \end{cases}
\end{equation}
For \(m=7\), every \(\kappa_p\) vanishes.  Pointwise
\(\sigma_p\ge0\) and \eqref{eq:C39-kappa-point-v2} make the two low points
have \(z=1\), while
the nine high points have \(z\in\{1,4\}\).  Since \(\sum z=35\), exactly
eight high points have \(z=4\), and the link equation makes their type
\((9,4,4)\).  Counting their incidences with five-blocks gives
\(8\cdot4\le2m=14\), contrary to
\cref{lem:n11-C39-interfaces-v2}.

For \(m=8\), all eleven points are high and precisely one has
\(\kappa=1\).  The ten zero-slack points have \(z\in\{1,4\}\), while the
exceptional point has \(z\in\{0,3\}\).  The sum \(\sum z=40\) therefore
forces respectively ten or nine points with \(z=4\).  Their incidences give
\(40\le16\) or \(36\le16\), again contradicting the harmonic cap.
\end{proof}

\begin{lemma}[The \(L=12\) low-host moment]
\label{lem:n11-C39-low-L12-v2}
In the \(C=39,L=12\) branch, some proper five-circle has host load
\(H\ge26\).
\end{lemma}

\begin{proof}
Suppose every proper five-circle has host load at most \(25\), and put
\[
 E=2m-9,\qquad z_p=d_5(p),\qquad
 G=2\sum_p\binom{z_p}{2}-\sum_{p\text{ high}}z_p.
\]
Here \(m\ge5\), \(\sum_pz_p=5m\), and the K1 moment is
\begin{equation}\label{eq:C39-global-moment-v2}
 \sum_\Delta H_\Delta=20m+G+\sum_\Delta q_\Delta.
\end{equation}
The successive costs of increasing \(z_p\) from zero to four are
\[
 \begin{array}{c|rrrr}
 &1&2&3&4\\ \hline
 p\text{ low}&0&2&4&6\\
 p\text{ high}&-1&1&3&5.
 \end{array}
\]
Filling these increasing marginal costs gives \(G\ge31\) at \(m=5\).
For \(m\ge6\), Cauchy and
\(\sum_{p\text{ high}}z_p\le4(2m-9)\) give
\[
 G\ge\frac{25m^2}{11}-13m+36>5m+5.
\]

The original-line Melchior row is
\(36+4\ell_4+9\ell_5\le52\), so there is at most one five-line.  A proper
five-circle contributes at most \(25\) to the left side of
\eqref{eq:C39-global-moment-v2}, and a possible five-line at most the
universal host cap \(30\).  Hence \(\sum H_\Delta\le25m+5\), whereas the
right side is at least \(131\) for \(m=5\) and is strictly larger than
\(25m+5\) for \(m\ge6\).  This contradiction proves the lemma.
\end{proof}

\begin{lemma}[The \(L=12\) maximum-host router]
\label{lem:n11-C39-maxhost-v2}
In the \(C=39,L=12\) branch, no proper five-circle of maximum host load has
\(26\le H\le30\).
\end{lemma}

\begin{proof}
Let \(h\) be the maximum proper-circle host load, let
\(t=\ell_5\in\{0,1\}\), and use \(m,E,G\) as in the preceding proof.
Besides \(G\ge31\) for \(m=5\), marginal filling gives
\begin{equation}\label{eq:C39-G-small-v2}
 G\ge45\quad(m=6),\qquad G\ge61\quad(m=7).
\end{equation}
and the Cauchy bound above is greater than \(8m+2\) for \(m\ge8\).
Since the possible five-line has host at most thirty, K1 gives
\begin{equation}\label{eq:C39-maxhost-moment-v2}
 20m+G\le h(m-t)+30t.
\end{equation}

For \(h=26\), \eqref{eq:C39-G-small-v2} and
\eqref{eq:C39-maxhost-moment-v2} leave
only \((m,t)=(5,1)\).  For \(h=27\), they
leave \((5,0),(5,1),(6,1)\).  But
\eqref{eq:C39-line-cut-v2} gives \(16t\le m+8\), removing every listed row
with \(t=1\).  Thus \(h\ne26\), and at \(h=27\) necessarily
\begin{equation}\label{eq:C39-H27-face-v2}
 m=5,\qquad t=0.
\end{equation}

We record a consequence used twice.  Rebase at a proper five-circle
\(\Delta\), put \(R=30-H_\Delta\), and let \(r,s,q,y\) have the front
meaning.  If \(R\ge3\) and \(s=0\), then
\begin{equation}\label{eq:C39-no-singleton-front-v2}
 A_{12}=3q-R,\qquad e_X=3-R-q+2y.
\end{equation}
If \(y=1\), then \(q=3\) and the second quantity is \(2-R<0\).  If
\(y=0\), nonnegativity in \eqref{eq:C39-no-singleton-front-v2} gives
simultaneously \(3q\ge R\) and
\(q\le3-R\), again impossible.  Hence every circle with \(H\le27\) has
\begin{equation}\label{eq:C39-singleton-degree-v2}
 s_\Delta\ge1.
\end{equation}

In the face \eqref{eq:C39-H27-face-v2}, form the graph whose five vertices
are the five-blocks and whose edges are their singleton intersections.  By
\eqref{eq:C39-singleton-degree-v2} it has minimum
degree one and hence at least three edges.  But
\eqref{eq:C39-singleton-v2} injects its edges into the \(E=1\) high points:
a carrier has \(d_5=2\) and can carry only that one edge.  This excludes
\(h=27\).

Suppose next that \(h=28\).
Equations \eqref{eq:C39-G-small-v2} and
\eqref{eq:C39-maxhost-moment-v2} leave
\(m\in\{5,6\}\),
and the line cut gives \(t=0\).  We claim that
\eqref{eq:C39-singleton-degree-v2} still holds for a circle
with \(H=28\).  Otherwise its outsider deficiency
\(\delta(e)=2-h_e\) has total weight two.  Its support is one zero edge
\(Z\), two adjacent single edges \(S_a\), or two disjoint single edges
\(S_d\).  Every \(A_{13}\) triangle contains a deficient edge by
\cref{lem:five-conic-rigidity}(K2.1).  A single edge has one host chord;
(K2.1) places that chord in a unique page-colour class, and triple ownership
allows at most one page of that colour through the edge.  Thus a zero edge
lies in four outsider triangles, while each single edge supports at most one
clean page, with the one
triangle on two adjacent singles as the only overlap.  Therefore
\begin{equation}\label{eq:C39-defect-page-caps-v2}
 r\le4,3,2\quad\hbox{in types }Z,S_a,S_d.
\end{equation}
With \(s=0\), the front says \(r+q=7\), so
\eqref{eq:C39-defect-page-caps-v2} and \(q\le3\) force type
\(Z\) and \((r,q)=(4,3)\).  If its zero edge is \(uv\), the four pages are
\(uvw\), one for every other outsider \(w\).  The high split gives
\(e_X=-2+2y\), hence \(y=1\): there is one outsider \(K_5\).  It must omit
\(u\) or \(v\), since otherwise it shares an outsider triple \(uvw\) with
a page, and it is consequently all-double.  Suppose it omits \(u\), and
let \(\gamma_i\) be the colour of the page \(uvw\).  The adjacent double
edges \(uw,vw\) determine two distinct vertices of \(D_i\); triple
ownership makes both host matchings omit \(i\).  The vertex coming from
\(vw\) lies on the all-double \(K_5\) trace and is its unique centre
\(\delta_i\) omitting \(i\), whereas the page trace contains
\(\gamma_i,\delta_i\) and the second vertex coming from \(uw\).  This is
exactly the tangent-separation obstruction in (K2.4).  The case in which
the \(K_5\) omits \(v\) is symmetric.  This proves the claim.

The singleton graph on the \(m\) proper five-circles consequently has
minimum degree one.  If \(m=5\), it has at least three singleton edges,
whereas \(E=1\), a contradiction.  Let \(m=6\).  Now it has at least three
edges and \eqref{eq:C39-singleton-v2} gives at most \(E=3\), so it is a
perfect matching and all three high points have \(z_p=2\).

Let \(A_j\) count pairs of five-circles meeting in \(j\) points, and put
\(Q=\sum_\Delta q_\Delta\).  Since \(A_1=3\),
\[
 A_0+A_1+A_2=15,\qquad
 \sum_p\binom{z_p}{2}=A_1+2A_2=27-2A_0.
\]
The three high points contribute six to the high incidence sum, hence
\begin{equation}\label{eq:C39-m6-G-v2}
 G=48-4A_0.
\end{equation}
Every disjoint pair contributes three to \(q_\Delta\) at each endpoint, so
\(Q\ge6A_0\).  The global moment and \(H_\Delta\le28\) give
\(G+Q\le48\).  Together with \eqref{eq:C39-m6-G-v2}, this forces
\begin{equation}\label{eq:C39-m6-equality-v2}
 A_0=0,\qquad Q=0,\qquad H_\Delta=28
 \quad\hbox{for every }\Delta.
\end{equation}
For any \(\Delta\), the perfect matching gives \(s=1\), while
\eqref{eq:C39-m6-equality-v2} gives
\(q=0\); the front therefore gives \(r=4\).  But an \(A_{14}\) outsider
four-set forces the two deficiency units to be the disjoint perfect
matching on that set.  Every \(A_{13}\) triangle then contains exactly one
of its two single edges, and (K2.1) permits at most one clean page on each.
Thus \(r\le2\), the final contradiction for \(h=28\).

If \(h=29\), the deficiency graph of a maximum-host circle is one single
edge \(uv\).  An \(A_{14}\) support is impossible, and every \(A_{13}\)
triangle contains \(uv\).  The edge has one host chord; (K2.1) assigns that
chord to a unique page colour, and triple ownership permits at most one page
of that colour through \(u,v\).  Hence \(r\le1\).  The front, however, gives
\(r+q=6\), so \(q\le3\) forces \(r\ge3\).

Finally, if \(h=30\), every outsider edge is double.  Clause (K2.1) forbids
an \(A_{13}\) block, and any \(A_{14}\) block contains a forbidden
all-double outsider triangle.  Thus \(r=s=0\), while the front gives
\(q=5>3\).  This exhausts \(26\le h\le30\).
\end{proof}
